% Section: Hyperparameter Optimization for Non-Stationary Routing
% Include from paper/sections/results.tex (main paper, not appendix)

\subsection{Hyperparameter Optimization for Non-Stationary Routing}
\label{sec:hparam_sweep}

\paragraph{The challenge.}
ParetoBandit's three core hyperparameters---exploration coefficient
$\alpha$, prior strength $\neff$, and forgetting factor
$\gamma$---interact non-trivially with the stationarity of the
deployment environment.
Standard cross-validation assumes i.i.d.\ data and optimises a
single stationary metric (e.g., reward or Pareto AUC), which
selects near-stationary forgetting ($\gamma \geq 0.999$)
because forgetting \emph{always} hurts stationary performance.
However, $\gamma{=}1.0$ fails catastrophically under distribution
shift: as shown in Section~\ref{sec:nonstationarity}, the
no-forgetting Naive Bandit accumulates $34\%$ more cumulative regret
than the tuned ParetoBandit ($\gamma{=}0.995$) after a reward swap.
Conversely, optimising solely for adaptability yields
$\gamma \ll 1$ with excessive forgetting that destabilises the
stationary policy.
This creates a fundamental tension absent from supervised routing:
\emph{the same hyperparameter must simultaneously protect stationary
efficiency and enable non-stationary recovery}.

Existing bandit-based LLM routers do not address this tension.
PILOT~\cite{panda2025pilot} grid-searches the exploration rate
$\alpha$ to maximise a single reward metric on held-out data,
with no forgetting mechanism and no non-stationarity consideration
in the selection protocol.
BaRP~\cite{wang2025barp} and LLM Bandit~\cite{li2025llmbandit}
omit forgetting entirely.
MixLLM~\cite{wang2025mixllm} delegates adaptation to periodic
neural-network retraining rather than tuning bandit-specific
parameters.
In the general bandit literature, the CDT framework~\cite{kang2024cdt}
proposes online hyperparameter tuning via a double-layer bandit that
treats $\alpha$ selection as a non-stationary continuum-armed problem,
extending earlier discrete methods~\cite{bouneffouf2020oplinucb}.
However, CDT remains \emph{single-objective} (maximise cumulative
reward) and does not tune the forgetting factor $\gamma$ or address
the efficiency--adaptability tradeoff.
The foundational Discounted UCB analysis by
Garivier and Moulines~\cite{garivier2011switching} derives $\gamma$
theoretically from the number of breakpoints, but this quantity is
unknown in every production deployment.
To the best of our knowledge, no prior work---in LLM routing or in
the broader non-stationary bandit literature---proposes a
multi-objective protocol for jointly selecting $(\alpha, \neff,
\gamma)$ to balance stationary efficiency against non-stationary
recovery.

\paragraph{Multi-objective formulation.}
We resolve this tension via an
$\epsilon$-constraint method~\cite{haimes1971bicriterion} that
treats hyperparameter selection as a bi-objective optimisation
problem over a grid of $(\alpha, \neff, \gamma)$ candidates
(grid: $\alpha \in \{0.01, 0.05, 0.1, 0.25, 0.5, 1.0\}$,
$\neff \in \{1, 10, 50, 200, 1{,}000, 5{,}000\}$,
$\gamma \in \{0.995, 0.997, 0.999, 1.0\}$):
%
\begin{enumerate}[nosep,leftmargin=*]
  \item \textbf{Primary objective} (efficiency):
    Maximise budget-paced Pareto AUC on the validation split, computed
    by sweeping log-spaced budget targets with the BudgetPacer
    active, building per-seed Pareto frontiers from the bandit's
    empirical routing points, and averaging AUC across 20~seeds.

  \item \textbf{Constraint}:
    Retain only configurations whose budget-paced AUC is within
    $\epsilon{=}0.25\%$ of the grid maximum, forming a \emph{feasible set}
    that protects stationary routing quality.

  \item \textbf{Secondary objective} (adaptability):
    Among feasible configurations, minimise the average Phase-2
    cumulative regret under a synthetic reward/cost swap
    (identical to the non-stationary protocol of
    Section~\ref{sec:reward_shift}, averaged over all
    $\binom{K}{2}$ arm-pair swaps and 20~seeds).
\end{enumerate}
%
This formulation gives budget-pacing \emph{priority}---no
configuration that materially degrades stationary routing can be
selected---while using the secondary objective to break ties among
near-optimal configurations by their ability to recover from
distribution shift.
The $\epsilon$ parameter gives the practitioner an explicit,
auditable knob: it specifies the maximum stationary quality they are
willing to sacrifice in exchange for non-stationary robustness.

\paragraph{Results.}
Table~\ref{tab:hparam_selection} summarises the selection outcome
for both ParetoBandit (warmup priors) and Tabula Rasa (cold start).

\begin{table}[t]
\centering
\caption{%
  \textbf{$\epsilon$-constraint hyperparameter selection.}
  AUC-only optimisation selects $\gamma{=}1.0$ (no forgetting)
  for both variants.
  For ParetoBandit, the $\epsilon$-constraint method
  ($\epsilon{=}0.25\%$) trades ${\sim}0.2\%$ AUC for moderate
  forgetting ($\gamma{=}0.997$) and stronger priors
  ($\neff{=}5{,}000$), lowering Phase-2 regret.
  For Tabula Rasa, the constraint \emph{binds}: forgetting
  imposes too steep an AUC penalty without priors, so the
  selected configuration coincides with the AUC-only winner.
  Feasible sets: 32/144 (ParetoBandit), 4/24 (Tabula Rasa).
  Metrics: budget-paced Pareto AUC (val, 20~seeds) and Phase-2
  cumulative regret (val, averaged over all arm-pair swaps,
  20~seeds).%
}
\label{tab:hparam_selection}
\small
\begin{tabular}{@{}llcccccc@{}}
\toprule
\textbf{Variant} & \textbf{Method}
  & $\boldsymbol{\alpha}$ & $\boldsymbol{\neff}$
  & $\boldsymbol{\gamma}$
  & \textbf{BP AUC} & \textbf{P2 Regret} \\
\midrule
\multirow{2}{*}{ParetoBandit}
  & AUC-only     & $0.01$ & $1$       & $1.000$ & $0.929$ & --- \\
  & $\epsilon$-constr.\ & $0.01$ & $5{,}000$ & $0.997$ & $0.926$ & $\mathbf{48.2}$ \\[3pt]
\multirow{2}{*}{Tabula Rasa}
  & AUC-only     & $0.10$ & ---    & $1.000$ & $0.926$ & --- \\
  & $\epsilon$-constr.\ & $0.10$ & ---    & $1.000$ & $0.926$ & $54.4$ \\
\bottomrule
\end{tabular}
\end{table}

AUC-only optimisation selects $\gamma{=}1.0$ (no forgetting) for
both variants.
Section~\ref{sec:nonstationarity} shows that no forgetting is
catastrophic under drift.
For ParetoBandit, the $\epsilon$-constraint method selects
$\gamma{=}0.997$ (effective memory ${\sim}333$ steps) with strong
priors ($\neff{=}5{,}000$), trading ${\sim}0.2\%$ budget-paced
AUC ($0.929 \to 0.926$) for substantially lower Phase-2 regret
($48.2$).
For Tabula Rasa, however, the constraint \emph{binds}: without
warmup priors, forgetting imposes a steep AUC penalty
(${\sim}0.9\%$ at $\gamma{=}0.995$) that exceeds the
$\epsilon{=}0.25\%$ tolerance, leaving only 4 of 24
configurations feasible---all with $\gamma \geq 0.997$.
The AUC-only winner ($\alpha{=}0.10$, $\gamma{=}1.0$) happens
to have the lowest regret among those four, so the
$\epsilon$-constraint method returns the same configuration.
On the held-out test split the selected configurations achieve
Pareto AUC of $\hpParetoBanditTestAUC$ (ParetoBandit) and
$\hpTabulaTestAUC$ (Tabula Rasa) versus $\hpFixedTestAUC$ for the
fixed-model envelope---deltas of $\hpParetoBanditTestDelta\%$ and
$\hpTabulaTestDelta\%$, both within one standard deviation of
seed variance ($\pm\hpParetoBanditTestStd$ and
$\pm\hpTabulaTestStd$).

Two patterns merit attention:

\begin{enumerate}[nosep,leftmargin=*]
  \item \textbf{Warmup priors are a prerequisite for the
    adaptability trade-off.}
    Only ParetoBandit can afford forgetting: its warmup priors
    provide a high-quality initial policy that buffers the AUC loss
    from discounting stale evidence, keeping the AUC within the
    $\epsilon$ tolerance.
    Without priors, Tabula Rasa must learn a 25-dimensional policy
    from scratch under budget pacing; any forgetting discards
    hard-won observations and precipitates a disproportionate AUC
    drop.
    This asymmetry demonstrates that the supervised initialization
    (Section~\ref{sec:warmup}) is not merely a convenience but a
    \emph{structural enabler} of the non-stationary adaptation
    mechanism.

  \item \textbf{The epsilon-constraint method shifts prior strength,
    not exploration.}
    For ParetoBandit, AUC-only selects $\neff{=}1$ (minimal priors)
    with $\gamma{=}1.0$; the $\epsilon$-constraint method raises
    $\neff$ to $5{,}000$ (the grid maximum) while introducing
    forgetting ($\gamma{=}0.997$).
    Stronger priors compensate for the AUC cost of forgetting,
    enabling the router to ``afford'' discounting stale evidence.
    Exploration ($\alpha{=}0.01$) remains unchanged,
    suggesting that the prior/forgetting axis---not the
    exploration axis---is the operative lever for balancing
    efficiency and adaptability in the warm-start regime.
\end{enumerate}
%
\noindent
Per-seed variance breakdowns and exhaustive per-configuration logs
are available in the accompanying results files.

\paragraph{Novelty and significance.}
This $\epsilon$-constraint protocol is, to our knowledge, the first
multi-objective hyperparameter selection framework for non-stationary
bandit routing.
It addresses a gap that spans both the LLM routing and the general
bandit literatures:
existing systems either assume stationarity in their tuning
protocol~\cite{panda2025pilot,wang2025barp,li2025llmbandit},
delegate adaptation to model retraining rather than tuning bandit
parameters~\cite{wang2025mixllm}, or derive $\gamma$ from
theoretical quantities (number of breakpoints) that are unknown in
production~\cite{garivier2011switching}.
The CDT framework~\cite{kang2024cdt} demonstrates that online
hyperparameter tuning is feasible for contextual bandits, but it
optimises a single reward objective and does not address the
stability--plasticity tradeoff that arises when the router must
simultaneously serve stationary traffic well and recover from shifts.

Our contribution is complementary: rather than tuning $\alpha$
online during execution (as CDT does), we propose a
\emph{deployment-time} protocol that jointly selects
$(\alpha, \neff, \gamma)$ using a multi-objective criterion
that explicitly encodes the dual requirements of
production routing---efficiency under stationarity and resilience
under drift.
The empirical result that single-objective tuning
selects $\gamma{=}1.0$ (no forgetting) for \emph{all} variants
is itself a cautionary finding: any practitioner who tunes via
standard cross-validation will converge on no forgetting,
offering zero protection against distribution shift---despite
this being the failure mode they are most likely to encounter
in production.
More strikingly, for cold-start variants the
$\epsilon$-constraint method \emph{also} selects $\gamma{=}1.0$
because forgetting is too costly without prior knowledge.
This reveals a fundamental asymmetry: the multi-objective
protocol can only unlock non-stationary robustness when warmup
priors cushion the stationary cost of forgetting.

\paragraph{Practical implications.}
From a deployment perspective, the protocol eliminates the need to
hand-tune forgetting factors on live traffic---a practice that risks
degrading routing quality during the tuning period and offers no
reproducibility.
The practitioner specifies two quantities: (i)~a grid of candidate
hyperparameters, and (ii)~the tolerance $\epsilon$ (how much
stationary quality they are willing to sacrifice for robustness).
The output is a single configuration that is
near-optimal for budget-paced routing and maximally adaptive within
that quality budget.
Because the secondary objective uses a synthetic reward swap
(Section~\ref{sec:reward_shift}), the protocol requires no live
non-stationary data and can be run entirely on a static validation
split before deployment.
This makes the tuning procedure safe, offline, and repeatable---a
necessary property for production systems subject to change-management
governance.
